%% kapa: exam=1 round=34 year=2023 period=2 form=A subject=appraisal_law subject_ko=감정평가관계법규 seq=40 number=40
%% source: Q-Net 감정평가사 기출문제 (cst00309) artlSeq=5211789 file=raw/1008/5211789/2023년도 제34회 감정평가사 1차시험 2교시.pdf
%% conversion: kordoc 4.24.51 -> tools/gen1.py (auto); review=proofread-sonnet; flags=none
\begin{question}{40}
동산·채권 등의 담보에 관한 법령상 동산담보권에 관한 설명으로 옳은 것은?
\begin{choices}
  \item 창고증권이 작성된 동산을 목적으로 담보등기를 할 수 있다.
  \item 담보권설정자의 사업자등록이 말소된 경우 그에 따라 이미 설정된 동산담보권도 소멸한다.
  \item 담보권설정자의 행위에 의한 사유로 담보목적물의 가액(價額)이 현저히 증가된 경우 담보목적물의 일부를 반환받을 수 있다.
  \item 동산담보권의 효력은 담보목적물에 대한 압류가 있은 후에 담보권설정자가 그 담보목적물로부터 수취할 수 있는 과실(果實)에 미친다.
  \item 담보권자와 담보권설정자 간 약정에 따른 동산담보권의 득실변경(得失變更)은 담보등기부에 등기하지 않더라도 그 효력이 생긴다.
\end{choices}
\end{question}
\answer{4}
